\documentclass[12pt,a4paper]{article}

% --- Pacotes Básicos ---
\usepackage[utf8]{inputenc}
\usepackage[T1]{fontenc}
\usepackage[brazil]{babel}
\usepackage{geometry}
\geometry{a4paper, left=3cm, right=2cm, top=3cm, bottom=2cm}
\usepackage{graphicx}
\usepackage{setspace}
\usepackage{hyperref}
\usepackage{xcolor}

% --- Configurações de Código ---
\usepackage{listings}
\definecolor{codegreen}{rgb}{0,0.6,0}
\definecolor{codegray}{rgb}{0.5,0.5,0.5}
\definecolor{codepurple}{rgb}{0.58,0,0.82}
\definecolor{backcolour}{rgb}{0.95,0.95,0.92}

\lstset{
    backgroundcolor=\color{backcolour},
    commentstyle=\color{codegreen},
    keywordstyle=\color{magenta},
    numberstyle=\tiny\color{codegray},
    stringstyle=\color{codepurple},
    basicstyle=\ttfamily\footnotesize,
    breakatwhitespace=false,
    breaklines=true,
    captionpos=b,
    keepspaces=true,
    numbers=left,
    numbersep=5pt,
    showspaces=false,
    showstringspaces=false,
    showtabs=false,
    tabsize=2
}

% --- Informações do Documento ---
\title{
    \large \textbf{Relatório de Desenvolvimento Didático: KeepCalm Web Browser} \\
    \small Projeto Acadêmico para Demonstração de Habilidades em Programação Full-Stack
}
\author{Estudante: Ryhan Schutz \\ Instituição: SENAI}
\date{\today}

\begin{document}

\maketitle
\onehalfspacing

\begin{abstract}
Este relatório documenta o processo de desenvolvimento do navegador KeepCalm Web, construído utilizando a infraestrutura Tauri v2, Rust e React. O projeto possui fins puramente didáticos, visando o aprimoramento de competências em desenvolvimento de sistemas, integração nativa e segurança da informação.
\end{abstract}

\section{Introdução}
O desenvolvimento de um navegador web contemporâneo apresenta desafios complexos que transcendem a simples exibição de páginas. Este projeto foi concebido como uma ferramenta de aprendizado prático no contexto das disciplinas de programação do SENAI. O objetivo principal não é competir comercialmente, mas sim compreender as nuances de comunicação entre processos (IPC), gerenciamento de estados no Frontend e segurança de memória no Backend usando a linguagem Rust.

\section{Arquitetura e Tecnologias}
Para a construção do software, foram selecionadas tecnologias de ponta que exigem alto rigor técnico:

\begin{itemize}
    \item \textbf{Tauri v2 (Framework)}: Utilizado para criar a ponte entre a interface web e o sistema operacional.
    \item \textbf{Rust (Backend)}: Responsável pelo motor de filtragem de rede e lógica de privacidade, garantindo performance e segurança.
    \item \textbf{React \& TypeScript (Frontend)}: Utilizados para a interface de usuário (UI), permitindo uma experiência dinâmica e tipagem forte.
    \item \textbf{Zustand}: Gerenciamento de estado leve para controle de abas e sessões.
\end{itemize}

\section{Desafios Técnicos e Aprendizado}
Durante a execução, diversos obstáculos surgiram, servindo como oportunidades pedagógicas cruciais.

\subsection{Gerenciamento de Client-Side Decorations (CSD)}
Um dos maiores desafios foi a implementação de uma barra de título personalizada que permitisse o arraste da janela nativa. Inicialmente, conflitos de eventos de mouse entre o React e o Tauri impediam o movimento suave da janela. A solução envolveu o estudo aprofundado de propagação de eventos e o uso de APIs nativas como \texttt{startDragging()}.

\subsection{Integração de WebViews Dinâmicas}
O gerenciamento de múltiplas abas exigiu uma sincronização precisa entre o estado do Zustand e o ciclo de vida das WebViews no Rust. Erros de \textit{lifetime} e concorrência no Rust foram diagnosticados e corrigidos, reforçando o aprendizado sobre o gerenciador de memória do Rust (Borrow Checker).

\section{Relatório de Erros e Diagnóstico}
Conforme a intenção de reportar o progresso ao SENAI, mantemos um log de erros que ilustra a evolução do projeto:

\begin{enumerate}
    \item \textbf{Erro de Permissões (Capabilities)}: Inicialmente as funções de minimizar/maximizar foram bloqueadas por segurança. \textbf{Solução}: Configuração correta do arquivo \texttt{default.json} no Tauri v2.
    \item \textbf{Conflito de UI}: Elementos do estilo \textit{Nocturnal Safari} estavam sobrepondo as áreas clicáveis. \textbf{Solução}: Ajuste de Z-Index e isolamento de eventos via \texttt{e.stopPropagation()}.
\end{enumerate}

\section{Conclusão}
O KeepCalm Web atingiu seus objetivos pedagógicos iniciais, apresentando um navegador funcional, seguro e com interface moderna. O processo demonstrou que a depuração constante e a documentação de erros são partes integrantes e valiosas da formação de um desenvolvedor de sistemas.

\end{document}
